\newpage
\section{Cześć teoretyczna}
Obszary do opisania, baza do wykorzystania w projekcie, dodatkowe informacje z tematu:

\subsection{Twelve Factor App}
\begin{enumerate}
    \item Kod źródłowy - Jedno źródło kodu śledzone systemem kontroli wersji, wiele wdrożeń
    \item Zależności - Jawnie zadeklaruj i wydziel zależności
    \item Konfiguracja - Przechowuj konfigurację w środowisku
    \item Usługi wspierające - Traktuj usługi wspierające jako przydzielone zasoby
    \item Buduj, publikuj, uruchamiaj - Oddzielaj etap budowania od uruchamiania
    \item Procesy - Uruchamiaj aplikację jako jeden lub więcej bezstanowych procesów
    \item Przydzielanie portów - Udostępniaj usługi przez przydzielanie portów
    \item Współbieżność - Skaluj przez odpowiednio dobrane procesy
    \item Zbywalność - Zwiększ elastyczność pozwalając na szybkie uruchamianie i zatrzymywanie aplikacji
    \item Jednolitość środowisk - Utrzymuj konfigurację środowisk jak najbardziej zbliżoną do siebie
    \item Logi - Traktuj logi jako strumień zdarzeń
    \item Zarządzanie aplikacją - Uruchamiaj zadania administracyjne jako jednorazowe procesy
\end{enumerate}

\subsection{DevOps}

\subsubsection{Kultura}
\begin{enumerate}
    \item \textbf{Współpraca i Komunikacja}: Promowanie otwartej i ciągłej komunikacji między zespołami w celu szybszego rozwiązywania problemów i lepszego zrozumienia celów.
    \item \textbf{Wspólna Odpowiedzialność}: Deweloperzy i specjaliści od operacji dzielą odpowiedzialność za cały cykl życia aplikacji, co zwiększa zaangażowanie i jakość pracy.
    \item \textbf{Ciągłe Doskonalenie}: Stałe dążenie do ulepszania procesów, narzędzi i produktów poprzez iteracyjne podejście i feedback.
    \item \textbf{Automatyzacja}: Wykorzystanie narzędzi do automatyzacji zadań takich jak testowanie, wdrażanie i monitorowanie w celu zwiększenia efektywności i redukcji błędów.
    \item \textbf{Uczenie się na Błędach}: Kultura, która traktuje błędy jako okazje do nauki, zachęcając do otwartości i dzielenia się wiedzą.
    \item \textbf{Pomiar i Analiza}: Używanie metryk i danych do podejmowania świadomych decyzji i monitorowania postępów.
    \item \textbf{Orientacja na Klienta}: Skupienie na dostarczaniu wartości dla użytkownika końcowego poprzez szybkie i niezawodne wydawanie oprogramowania.
    \item \textbf{Eliminacja Silosów}: Rozbijanie barier między działami w organizacji w celu tworzenia zintegrowanego i spójnego środowiska pracy.
\end{enumerate}

\subsubsection{Cykl Ciągłego Doskonalenia (\textit{ang.} Continuous Improvement)}
Kluczowe (kolejne podrozdziały) do realizacji części praktycznej
\begin{enumerate}
    \item \textbf{Planowanie} z wykorzystaniem narzędzi do zarządzania projektami.
    \item \textbf{Kodowanie} przy użyciu nowoczesnych technologii i narzędzi programistycznych.
    \item \textbf{Budowanie} i \textbf{testowanie} w zautomatyzowanych pipeline'ach CI/CD.
    \item \textbf{Wydawanie} stabilnych wersji aplikacji poprzez konteneryzację i zarządzanie pakietami.
    \item \textbf{Wdrażanie} na skalowalnej infrastrukturze chmurowej z wykorzystaniem IaC i orkiestracji kontenerów.
    \item \textbf{Operacje} z zapewnieniem wysokiej dostępności i elastyczności dzięki usługom chmurowym.
    \item \textbf{Monitorowanie} aplikacji i infrastruktury w celu zapewnienia jakości i szybkiego reagowania na incydenty.
\end{enumerate}

\subsubsection{Pochodne od DevOps, nowe kierunki, specjalizacje}
\begin{itemize}
% Wykorzystane w części praktycznej:
    \item \textbf{DevSecOps}: Integracja bezpieczeństwa (Security) w procesie DevOps. Bezpieczeństwo jest uwzględniane na każdym etapie cyklu życia aplikacji, od planowania po wdrożenie i utrzymanie, co pozwala na szybsze wykrywanie i eliminowanie zagrożeń.
    \item \textbf{GitOps}: Podejście, które wykorzystuje Git jako jedyne źródło prawdy dla infrastruktury i aplikacji. Automatyzacja wdrożeń jest osiągana poprzez monitorowanie zmian w repozytorium Git i automatyczne zastosowanie ich w środowisku produkcyjnym.
% Specjalizacje inżynierskie
    \item \textbf{Cloud Engineering}:
    \item \textbf{Site Reliability Engineering}:
    \item \textbf{Platform Engineering}:
% Pozostałe:    
    \item \textbf{FinOps}: Połączenie praktyk finansowych z operacjami IT. Celem jest optymalizacja kosztów chmury poprzez lepszą współpracę między działami finansowymi, technicznymi i biznesowymi, co prowadzi do świadomego zarządzania wydatkami na infrastrukturę.
    \item \textbf{MLOps}: Integracja praktyk DevOps w dziedzinie uczenia maszynowego (Machine Learning). Ułatwia zarządzanie cyklem życia modeli ML, od trenowania po wdrażanie i monitorowanie, zapewniając spójność i skalowalność modeli w produkcji.
    \item \textbf{DataOps}: Zastosowanie zasad DevOps w inżynierii danych. Skupia się na automatyzacji procesów związanych z danymi, poprawie ich jakości oraz współpracy między zespołami, co przyspiesza dostarczanie wartościowych informacji biznesowych.
    \item \textbf{AIOps}: Wykorzystanie sztucznej inteligencji (AI) w operacjach IT. Pomaga w automatyzacji i optymalizacji zarządzania infrastrukturą poprzez analizę dużych zbiorów danych i predykcyjne wykrywanie problemów.
    \item \textbf{ModelOps}: Skupia się na zarządzaniu cyklem życia modeli analitycznych i AI w produkcji. Ułatwia wdrażanie, monitorowanie i aktualizację modeli w środowiskach produkcyjnych.
    \item \textbf{NoOps}: Koncepcja zakładająca pełną automatyzację operacji IT, eliminując potrzebę tradycyjnych działań operacyjnych. Deweloperzy mogą skupić się wyłącznie na kodzie, podczas gdy infrastruktura zarządza się sama poprzez chmurę i automatyzację.
    \item \textbf{BizDevOps}: Rozszerzenie DevOps o aspekt biznesowy. Ma na celu zintegrowanie celów biznesowych z procesem deweloperskim i operacyjnym, co prowadzi do lepszego dopasowania produktów do potrzeb rynku.
    \item \textbf{SecOps}: Skupia się na integracji bezpieczeństwa z operacjami IT, często w czasie rzeczywistym. Chociaż podobne do DevSecOps, SecOps kładzie większy nacisk na ciągłe monitorowanie i reagowanie na incydenty bezpieczeństwa.
    \item \textbf{ChatOps}: Wykorzystanie platform komunikacyjnych (np. Slack, Microsoft Teams) do zarządzania operacjami i automatyzacją zadań poprzez interakcje w czasie rzeczywistym, co zwiększa przejrzystość i szybkość reakcji zespołów.
    \item \textbf{IoTOps}: Integracja zasad DevOps w obszarze Internetu Rzeczy (Internet of Things). Pomaga w zarządzaniu i monitorowaniu urządzeń IoT oraz zapewnia bezpieczne i efektywne wdrażanie aktualizacji.
\end{itemize}

\subsection{Cloud computing}

\subsubsection{Rodzaje chmur}
\begin{enumerate}
    \item \textbf{Chmura publiczna} \\
    Chmura publiczna to model usług chmurowych, w którym zasoby obliczeniowe, takie jak serwery, przestrzeń dyskowa, aplikacje i moc obliczeniowa, są dostarczane przez zewnętrznego dostawcę usług za pośrednictwem internetu. W tym modelu infrastruktura jest współdzielona między wieloma użytkownikami (klientami), a dostawca chmury jest odpowiedzialny za zarządzanie, utrzymanie i zabezpieczenie zasobów. Użytkownicy płacą zazwyczaj za rzeczywiste wykorzystanie zasobów według modelu "pay-as-you-go". Chmury publiczne oferują wysoką skalowalność i elastyczność, co pozwala na szybkie dostosowanie się do zmieniających się potrzeb biznesowych bez konieczności inwestowania w infrastrukturę sprzętową. Przykładami dostawców chmury publicznej są:
    \begin{itemize}
        \item \textbf{Amazon Web Services (AWS)}
        \item \textbf{Microsoft Azure}
        \item \textbf{Google Cloud Platform (GCP)}
    \end{itemize}

    \item \textbf{Chmura prywatna} \\
    Chmura prywatna to infrastruktura chmurowa wykorzystywana wyłącznie przez jedną organizację. Może być zarządzana przez tę organizację lub przez zewnętrznego dostawcę usług, ale zasoby i usługi są dostępne tylko dla tej konkretnej firmy, co zapewnia wyższy poziom kontroli nad danymi, bezpieczeństwem i zgodnością z regulacjami. Chmury prywatne są często wykorzystywane przez przedsiębiorstwa, które mają ścisłe wymagania dotyczące bezpieczeństwa, poufności danych lub regulacji prawnych (np. instytucje finansowe, rządowe, medyczne). Mogą być hostowane w lokalnym centrum danych organizacji lub w dedykowanym ośrodku danych dostawcy usług. Chmura prywatna oferuje:

    \begin{itemize}
        \item \textbf{Wyższy poziom bezpieczeństwa i prywatności}: Dane są przechowywane na dedykowanych serwerach.
        \item \textbf{Dostosowanie}: Możliwość konfiguracji infrastruktury zgodnie z konkretnymi potrzebami.
        \item \textbf{Kontrola}: Pełny wgląd i zarządzanie zasobami oraz politykami bezpieczeństwa.
    \end{itemize}

    \item \textbf{Chmura hybrydowa} \\
    Chmura hybrydowa to połączenie chmury publicznej i prywatnej, które pozwala na przenoszenie danych i aplikacji między tymi dwoma środowiskami. Umożliwia to organizacjom korzystanie z zalet obu modeli:

    \begin{itemize}
        \item \textbf{Elastyczność i skalowalność chmury publicznej}: Możliwość szybkiego zwiększenia zasobów w razie potrzeby.
        \item \textbf{Bezpieczeństwo i kontrola chmury prywatnej}: Przechowywanie wrażliwych danych w zabezpieczonym środowisku.
    \end{itemize}

    Chmura hybrydowa jest idealna dla firm, które chcą utrzymać kontrolę nad krytycznymi danymi i aplikacjami, jednocześnie korzystając z mocy obliczeniowej chmury publicznej do obsługi zmiennego obciążenia lub wykonywania zadań wymagających dużej mocy obliczeniowej. Przykładowe scenariusze użycia chmury hybrydowej:

    \begin{itemize}
        \item \textbf{Disaster Recovery}: Kopie zapasowe danych w chmurze publicznej.
        \item \textbf{Burst Computing}: Przenoszenie dodatkowego obciążenia do chmury publicznej podczas szczytów.
        \item \textbf{Testowanie i rozwój}: Wykorzystanie chmury publicznej do tworzenia i testowania aplikacji przed wdrożeniem w środowisku prywatnym.
    \end{itemize}

    \item \textbf{Chmury specjalne (branżowe, społecznościowe)} \\
    Chmury specjalne, znane również jako chmury społecznościowe lub branżowe, są zaprojektowane w celu spełnienia specyficznych wymagań określonej grupy użytkowników, branży lub społeczności. Infrastruktura jest współdzielona między organizacjami o podobnych potrzebach, takich jak:

    \begin{itemize}
        \item \textbf{Regulacje prawne i zgodność}: Branże podlegające restrykcyjnym regulacjom (np. służba zdrowia, finanse).
        \item \textbf{Wymagania dotyczące bezpieczeństwa}: Potrzeba wysokiego poziomu ochrony danych.
        \item \textbf{Specyficzne potrzeby branżowe}: Wymagania dotyczące wydajności, integracji czy specjalistycznych aplikacji.
    \end{itemize}

    Chmury specjalne łączą zalety chmur prywatnych (bezpieczeństwo, kontrola) i publicznych (skalowalność, efektywność kosztowa) w kontekście specyficznych potrzeb. Przykłady chmur specjalnych:

    \begin{itemize}
        \item \textbf{Chmura dla sektora publicznego}: Spełniająca wymogi rządowe dotyczące bezpieczeństwa i przechowywania danych.
        \item \textbf{Chmura dla służby zdrowia}: Zgodna z regulacjami dotyczącymi ochrony danych pacjentów (np. RODO, HIPAA).
        \item \textbf{Chmura edukacyjna}: Wspierająca instytucje edukacyjne w dostarczaniu zasobów i aplikacji dla studentów i pracowników.
    \end{itemize}

\end{enumerate}

\subsubsection{Usługodawcy do wyboru}
\begin{itemize}
    \item AWS
    \item Azure
    \item GCP
\end{itemize}

\subsubsection{Typy usług}
\begin{itemize}
    \item SaaS (Software as a Service): Oprogramowanie jako usługa. Dostarczanie aplikacji przez internet, dostępnych za pośrednictwem przeglądarki, bez konieczności instalacji lokalnej. Przykłady: Google Workspace, Microsoft Office 365, Salesforce.
    \item PaaS (Platform as a Service): Platforma jako usługa. Udostępnianie środowiska do tworzenia, testowania i wdrażania aplikacji bez zarządzania infrastrukturą. Przykłady: Google App Engine, Microsoft Azure App Service, Heroku.
    \item IaaS (Infrastructure as a Service): Infrastruktura jako usługa. Dostarczanie zasobów obliczeniowych, takich jak serwery, sieci, przestrzeń dyskowa, na żądanie. Przykłady: Amazon Web Services (AWS), Microsoft Azure, Google Compute Engine.
    \item FaaS (Function as a Service): Funkcja jako usługa. Model serverless, w którym kod jest uruchamiany w odpowiedzi na zdarzenia, a infrastruktura jest zarządzana przez dostawcę. Przykłady: AWS Lambda, Azure Functions, Google Cloud Functions.
    \item CaaS (Container as a Service): Kontenery jako usługa. Usługi zarządzania kontenerami, ułatwiające wdrażanie i skalowanie aplikacji w kontenerach. Przykłady: AWS Elastic Container Service, Google Kubernetes Engine, Azure Kubernetes Service.
    \item BaaS (Backend as a Service): Backend jako usługa. Udostępnianie zaplecza dla aplikacji mobilnych i webowych, takich jak bazy danych, uwierzytelnianie, przechowywanie plików. Przykłady: Google Firebase, AWS Amplify, Backendless.
    \item XaaS (Anything as a Service): Wszystko jako usługa. Ogólny termin obejmujący wszystkie usługi dostarczane przez chmurę, umożliwiający elastyczne skorzystanie z usługi, zasobów IT i ekspertyzy dostawcy z danego obszaru (przykłady XaaS: Gaming, Learning, Voice, Knowledge, Mobility, Logging, Encryption, etc)
\end{itemize}

\subsection{Cloud Native}
\begin{itemize}
    \item co to jest? różne definicje
    \item Cloud computing vs cloud native - implementacja vs architektura
    \item Cloud-native pillars
    \item jaki jest wplyw na IT? Cloud-native benefits
\end{itemize}

\subsubsection{Cloud Native Computing Fundation}
\begin{itemize}
    \item technologie i narzedzia - Landscape
    \item Standardy
    \item Eventy
    \item Kursy i certyfikaty
\end{itemize}

\subsection{Kubernetes}
zloty standard, no1 projekt CNCF, dominacja rynku
\begin{itemize}
    \item główne zalety
        \begin{itemize}
            \item everything as code
            \item cloud agnostic
            \item spójność narzędzi niezależnie od dostawcy
            \item automatyzacja
            \item rozszerzalne API
            \item deklaratywnosc i self healing cycle
        \end{itemize}
    \item alternatywy:
        \begin{itemize}
            \item rancher
            \item openshift
        \end{itemize}
    \item managed kubernetes:
        \begin{itemize}
            \item GKE
            \item EKS
            \item AKS
        \end{itemize}
    \item cloud containers:
        \begin{itemize}
            \item ECS
            \item Azure Container Instances
            \item Google Cloud Run
        \end{itemize}
    \item lekka orkiestracja:
        \begin{itemize}
            \item Docker Swarm
            \item Nomad
        \end{itemize}
\end{itemize}





